% !Mode:: "TeX:UTF-8"
% !TEX root = ..\main.tex
\chapter{绪论\footnote{章标题中脚注命令测试}(or 引言)各种测试}
绪论占坑，但是也要测试到底占多少缩进，换行情况，行距，都是这些不听话的小伙伴，好好调教你们。

\section{列表环境\footnote{节标题中脚注命令测试}测试}
以下是一个测试用的列表环境，内容不要在意。\footnote{正文中中脚注命令测试}

这里测试列表标签功能的交叉引用格式\ref{itm:11},\ref{itm:12},\ref{itm:13},\ref{itm:14}，分别表示第一至第四层级的itemize系列的交叉引用情况。
\begin{enumerate}
	\item 第一级列表\label{itm:11}
	\item 第一级列表
	\begin{enumerate}
		\item 第二级列表\label{itm:12}
		\item 第二级列表
		\begin{enumerate}
			\item 第三级列表\label{itm:13}
			\item 第三级列表
			\begin{enumerate}
				\item 第四级列表\label{itm:14}
				\item 第四级列表
				\item 第四级列表
				\item 第四级列表
			\end{enumerate}
			\item 第三级列表
			\item 第三级列表
			\item 第三级列表
		\end{enumerate}
		\item 第二级列表
		\item 第二级列表
	\end{enumerate}
	\item 第一级列表
	\item 第一级列表
	\item 第一级列表
\end{enumerate}

正是由于油膜物质的发现，使“雾伞”计划成为可能，这个计划是用核爆炸在太空中蒸发和扩散油膜物质，在太阳与地球之间形成一团“油膜尘埃”，降低太阳 对地球的辐射，达到缓解地球温室效应的目的。“我记得，海王星轨道附近应该还有前战争时期的恒星型核弹吧？”肯又问。“有的，‘雾伞’工程的飞船也装载了一些，在海王星环和卫星上爆破用，具体数目不清楚。” “好像一颗就够了。”肯兴奋起来。两个世纪前面壁者雷迪亚兹的战略计划中所研制的恒星型氢弹，后来共制造了五千多颗。虽然这种武器在末日之战中作用有限，但正如雷迪亚兹所言，各大 国主要是为可能爆发的人类之间的行星际战争准备的，核弹主要在大低谷时期制造，那时由于资源的匮乏，国际关系极其紧张，人类自身的战争一触即发。进入新时期后，这些骇人听闻的武器成了危险的鸡肋，虽然其所有权都属于地球国家， 但还是都被送入太空存贮，少部分已经用于行星工程的爆破，还有一部分送入太阳系外围轨道。曾有人设想将核弹中的聚变材料可以作为远程飞船的燃料补充，但由于核弹的拆解很困难，这个设想一直没有真正实现过\footnote{看看另起一页脚注编号的变化}。
\begin{itemize}
	\item 第一级列表
	\item 第一级列表
	\begin{itemize}
		\item 第二级列表
		\item 第二级列表
		\begin{itemize}
			\item 第三级列表
			\item 第三级列表
			\begin{itemize}
				\item 第四级列表
				\item 第四级列表
				\item 第四级列表
				\item 第四级列表
			\end{itemize}
			\item 第三级列表
			\item 第三级列表
			\item 第三级列表
		\end{itemize}
		\item 第二级列表
		\item 第二级列表
	\end{itemize}
	\item 第一级列表
	\item 第一级列表
	\item 第一级列表
\end{itemize}

“你觉得能行，”罗宾逊两眼放光地问道，他后悔这么简单的事自己怎么没 想到，一个载入史册的机会让肯抢去了。“试试吧，只有这一个办法了。”“如果行，博士，以后林格一斐兹罗监测站将永远按产生1G重力的速度旋转。”“这可是人类造出来的最大的东西了。”“蓝影”号飞船的指令长看着舱外漆黑的太空说，他极力想象自己能看到尘埃云，但确实什么
\begin{enumerate}
	\item 第四级列表
	\item 第四级列表
	\begin{enumerate}
		\item 第五级列表
		\item 第五级列表
		\item 第五级列表
	\end{enumerate}
	\item 第四级列表
	\item 第四级列表
\end{enumerate}
都看不到。“为什么它不能被阳光照出来呢，就像彗星的尾巴那样...”飞船驾驶员说，“蓝影”号上只有他和指令长两个人。他知道，尘埃云的密度确实像彗星尾一样稀薄，几乎和地球上实验室中造出的真空差不多。“可能是阳光太弱吧。”指令长回头看看太阳，在这海王星轨道和柯伊伯带 之间的冷寂空间，太阳看上去只是一颗刚能看出圆盘形状的大星星。阳光倒是还可以在舱壁上照出亮影，但已经十分微弱了。“再说，彗尾也要在一定的距离外 才能看到，我们可是就在云的边缘。”

\section{参考文献测试}
测试一下引用\cite{shi_chinas_2010}，引用\cite{shi2010china,hata2014soi,muhammad2011development}，还有其它引用\cite{shi2010china,muhammad2011development,lamport1994latex}.

\section{浮动体测试}
\subsection{插图测试}
如图\ref{fig:first_image_tset}是对此模版的第一张插图测试。

\begin{table}[htbp]
	\centering
	\includegraphics[width = 0.5\linewidth]{Chapter1.png}
	\caption{第一张插图测试}\label{fig:first_image_tset}
\end{table}

以下是一段对这些插图来历的介绍，引用自知乎专栏All about TeXnique中夏晓昊的文章\href{http://zhuanlan.zhihu.com/LaTeX/19669122}{《The TeXbook导读：从那头（多图杀猫的）狮子说起》}。

在The TeXbook中，有着一系列的以狮子为主题的插图。这些插图的作者是Duane Bibby。也是从The TeXbook开始，不少TeX书也采取了以狮子为主的插图，作者也是Duane Bibby。另外，每年的TUG（TeX Users Group）年会都会有一张以狮子为主题的logo，这只狮子已经是社区的吉祥物了。

为什么选择狮子呢？Yannis Haralambous写道（原文法语，此为转译后的英文）：Not for nothing is TeX represented by a lion. Donald Knuth has told us that lions are to him the guardians of libraries in the United States because there is a statue of a lion in front of the entrance of each large library there. Guardian of libraries, guardian of the Book—is that not indeed what TeX ultimately aspires to be? 或许吧。 （顺便说一句，TeX和MetaFont都用了狮子，TeX是公狮子，MetaFont是母狮子，多么和谐的一对啊。如果你还是忽略MetaFont的存在，那你还没有认识到它的重要性。）

作为插图，首要的一点就是贴切，然后是有趣。在TeX社区里面，have fun是一个很重要的词组，也有人说Happy TeXing。我知道有不少人不喜欢TeX，但是能有什么理由呢？如果你用不到它，那么浅尝辄止即可。如果你会用到很频繁，最好慢慢修炼做到精通。如果你只是偶尔用到，那么可以搬个模版什么的，甚至也可以找人帮你（不要指望别人会用足够的空闲时间来帮你，他没有这个义务，请支付报酬，最少也得请吃个红烧肉吧）。下面的插图，是TeX TeXbook中的，我也希望这个新年的假期，能有人有空来看看这本书。即使不能把所有的东西都看懂，那么也会对TeX的设计有了一定的了解，拿到扳手就好。

